\subsection{Базовые задачи движения под углом к горизонту}
%%%%%%%%%%%%%%%%%%%%%%%%%%%%%%%%%%%%%%%%%%%%%%%%%%%%
%%%%%%%%%%%%%%%%%%%%%%%%%%%%%%%%%%%%%%%%%%%%%%%%%%%%
\z{Камень, брошенный под углом $\alpha = 30\degree$ к горизонту, дважды побывал на одной высоте $h$ спустя время $t_1 =3\, \text{c}$ и $t_2  = 5\, \text{c}$ после начала движения. Найдите начальную скорость камня $v_0$ и высоту $h$.}
%
{$v_0=\frac{g(t_1+t_2)}{2\sin{\alpha}}=78.4\, \nicefrac{\text{м}}{\text{с}}$;~$h=\frac{1}{2}gt_1t_2=73.5\,\text{м}$}

%%%%%%%%%%%%%%%%%%%%%%%%%%%%%%%%%%%%%%%%%%%%%%%%%%%%

 \z{Тело бросают с высоты $h = 4\, \text{м}$ вверх под углом $\alpha = 45\degree$ к горизонту так, что к поверхности земли оно подлетает под углом $\beta = 60\degree$. Какое расстояние по горизонтали пролетит тело?}
%
{$L=\frac{2h}{\tan{\beta}-\tan{\alpha}}=10.9\,\text{м}$}
	
%%%%%%%%%%%%%%%%%%%%%%%%%%%%%%%%%%%%%%%%%%%%%%%%%%%%
\z{Из одной точки горизонтально в противоположных направлениях одновременно вылетают две частицы с начальными скоростями $v_1$ и $v_2$. Через какое время угол между скоростями частиц станет равным $90\degree$? Ускорение свободного падения равно $g$.}
%
{$t=\frac{\sqrt{v_1v_2}}{g}$}

%%%%%%%%%%%%%%%%%%%%%%%%%%%%%%%%%%%%%%%%%%%%%%%%%%%%	
\z{С высоты $h$ на наклонную плоскость, образующую с горизонтом угол $\alpha$, свободно падает мяч и упруго отскакивает от нее. Через какое время $t$ после удара мяч снова упадет на наклонную плоскость? Найдите расстояние от места первого удара до второго, от второго до третьего и т.д.}
%
{$t=2\sqrt{\frac{2h}{g}}$;~$s_1=8h\sin{\alpha}$;~$s_N=s_1\cdot N$}

%%%%%%%%%%%%%%%%%%%%%%%%%%%%%%%%%%%%%%%%%%%%%%%%%%%%
\z {Мячик бросают с начальной скоростью $v$ с поверхности земли под углом $\alpha$ к горизонту. В момент нахождения мячика на максимальной высоте из той же точки на поверхности земли бросают камень под углом $\beta$ к горизонту. Размеры мячика и камня малы, сопротивлением воздуха можно пренебречь.%

	\lpu Определите, с какой начальной скоростью $u$ бросили камень, если он столкнулся с мячиком во время его полёта.\lpuend
	\lpu Найдите время движения камня от момента его броска до момента столкновения с мячиком. \lpuend
}
%
{\puan$u=\frac{v\cos{\alpha}}{\cos{\beta}-2\ctg{\alpha}\sin{\beta}}$;~\puan$t=\frac{v\sin^2{\alpha}(\cos{\beta}-2\ctg{\alpha}\sin{\beta})}{2g\cos{\alpha}\sin{\beta}}$}

